\chapter{Mevcut Durum ve Sonuçlar}

\section{Toplam Test Kapsamı}
Proje $\sim$273 data-driven senaryo içermektedir. UI regression suite (\texttt{ui-regression.xml}) 14 test sınıfını 6 TestNG grubunda koşturur.

\section{Modül Bazlı Sonuç Özeti}
% TODO: Son mvn clean test çıktısından güncel pass/fail sayılarını tabloya yaz.

\begin{table}[H]
  \centering
  \caption{Modül bazlı senaryo dağılımı (hedef)}
  \begin{tabular}{@{}lrr@{}}
    \toprule
    \textbf{Modül} & \textbf{Senaryo} & \textbf{Risk} \\
    \midrule
    Checkout & 28 & Yüksek \\
    Login & 24 & Yüksek \\
    ForgotPassword & 24 & Yüksek \\
    Register & 22 & Orta-Yüksek \\
    Cart & 22 & Orta-Yüksek \\
    MyAccountSettings & 22 & Orta-Yüksek \\
    MyAccount & 20 & Orta \\
    Search & 18 & Orta \\
    WishList & 12 & Düşük \\
    Newsletter & 10 & Düşük \\
    Api & 32 & — \\
    Db & 35 & — \\
    E2E & 4 & Yüksek \\
    \bottomrule
  \end{tabular}
\end{table}

\section{Bilinçli Güvenlik Test Fail'leri}
ForgotPassword modülünde OWASP odaklı senaryolar (\texttt{FPW-EG-01}, \texttt{FPW-ST-02}) bilinçli olarak fail tasarlanmıştır. Bu senaryolar gerçek güvenlik açıklarını test etmek içindir; framework hatası değildir.

\section{Performans ve Paralel Koşum}
TestNG \texttt{parallel="classes"} modu ile 4 thread kullanılır. ThreadLocal WebDriver ile sınıf düzeyinde izolasyon sağlanır.

\section{Karşılaşılan Zorluklar ve Çözümler}
\begin{itemize}
  \item \textbf{DOM değişiklikleri:} Self-healing locator zinciri.
  \item \textbf{Dashboard veri kaybı:} Modül koşumlarında \texttt{mvn clean} kullanılmaması.
  \item \textbf{Ek dosya yolları:} \texttt{resolveDashboardAsset()} ile path normalizasyonu.
  \item \textbf{Ürün katalog değişimi:} JSON senaryo güncellemesi (ör. WishList ürün adı).
\end{itemize}

\section{Proje Değerlendirmesi}
Proje; üç katmanlı otomasyon, self-healing, AI analiz ve görsel dashboard ile bitirme projesi kapsamını karşılamaktadır. Eksik kalan CI/CD ve Jira REST API entegrasyonu gelecek çalışma olarak planlanmıştır.
